\documentclass[11pt,a4paper]{article}
\usepackage[utf8]{inputenc}
\usepackage[T1]{fontenc}
\usepackage[margin=1in]{geometry}
\usepackage{xcolor}
\usepackage{hyperref}
\usepackage{parskip}

\definecolor{primarycolor}{RGB}{33,37,41}
\definecolor{secondarycolor}{RGB}{0,102,204}
\hypersetup{colorlinks=true,urlcolor=secondarycolor,linkcolor=secondarycolor}
\pagestyle{empty}
\setlength{\parindent}{0pt}

\begin{document}

Dear Hiring Manager,

I am applying for the Cloud Infrastructure / DevOps Contractor role supporting Corvale Mobility’s fleet connectivity platform. In my current role, I own CI pipelines and build ephemeral test environments for three product teams, experience that aligns closely with the role’s focus on dependable infrastructure and services that developers can deploy and operate effectively.

At Odrava Health Software, I reduced average pipeline duration from 38 to 14 minutes by parallelising jobs, caching dependencies and moving end-to-end tests into on-demand ephemeral environments. I built those environments with Docker Compose and later Helm charts on a shared Kubernetes cluster. I also maintain a pytest and Playwright suite of approximately 900 tests, bringing the flaky-test rate from 7\% to under 1\%, and have written Python and Bash tooling for PostgreSQL test-data seeding. These responsibilities have involved making improvements across teams and keeping development and test workflows reliable.

My experience is strongest in CI/CD and test infrastructure rather than Terraform-based cloud infrastructure design. I would bring a practical foundation in Linux, containers, Helm, Kubernetes, scripting and operational reliability to Corvale’s work taking train-to-cloud connectivity from architecture toward production. The opportunity to support a platform involving telemetry, alerting and remote updates is a compelling application of that infrastructure experience.

Thank you for considering my application. I would welcome the opportunity to discuss how my experience improving pipelines, test environments and developer workflows could contribute to the fleet connectivity team.

Sincerely,

Marta Kowal-Brandt

\end{document}